% Chapter 2

\chapter{\uppercase{Literature Survey}}
\label{ch:survey}
This chapter reviews existing work on long-term memory for agents. All numbers are the ones the
authors reported; we did not rerun these systems.

\section{\uppercase{Long Context and External Memory}}
The simplest memory keeps the whole conversation in the prompt. Models use the middle of a long
context less reliably than its start or end \cite{liu2024lost}, and it is slow: the Mem0 paper
reports 17\,s p95 latency for the full history. The systems below keep memory outside the prompt
and show the model a selected part. They differ in what they store, whether old values are kept,
and whether the model or the code decides updates.

\section{\uppercase{Memory Systems}}
\textbf{MemGPT} \cite{packer2023memgpt} splits the prompt into a main context and external
storage, and the model moves data with function calls. It reports 92.5\% on deep memory
retrieval against 32.1\% for a summary. A missed function call loses the update.

\textbf{Zep and Graphiti} \cite{rasmussen2025zep} store a temporal knowledge graph. Each fact
has a valid and an invalid time, and a contradicting fact closes the old one without deleting
it. They report 71.2\% on LongMemEval against 60.2\% for the full history. This is closest to
our design, but Zep's main store is the graph, while ours is PostgreSQL.

\textbf{Mem0} \cite{chhikara2025mem0} has an LLM extract candidate facts, compare each with the
nearest stored memories and choose add, update, delete or no-op. On LoCoMo it reports 66.9
against 72.9 for the full history, at 1.4\,s p95 instead of 17\,s. We use the same
extract-then-compare idea, but fixed code makes the update decision.

\textbf{Supermemory} \cite{daga2026supermemory} links atomic memories to source chunks with
update, extend and derive links. Its results come only from a vendor report.
\textbf{Hindsight} \cite{latimer2025hindsight} builds on retain, recall and reflect; recall fuses
four retrieval strategies into a token budget. Our recall follows this design.

\textbf{Benchmarks.} LongMemEval \cite{wu2025longmemeval} and LoCoMo \cite{maharana2024locomo}
test long-term conversational memory. LoCoMo is the only benchmark in the Mem0 paper. Published
results are hard to compare because each paper uses its own models and judge.

\section{\uppercase{Methods Used in This Project}}
Keyword search uses PostgreSQL full-text search \cite{postgres2026fts} with GIN indexes. pgvector
\cite{kane2026pgvector} adds 1,024-dimensional vectors with cosine distance, produced by a Qwen3
embedding model \cite{zhang2025qwen3emb}. Reciprocal Rank Fusion \cite{cormack2009rrf} merges
ranked lists by summing $1/(60+\mathit{rank})$. The memory model is fine-tuned with QLoRA
\cite{dettmers2023qlora}, which trains small adapters on a frozen, quantised base model.

\section{\uppercase{Comparison and Gap}}
\begin{table}[H]
\caption{Comparison of existing memory systems}
\label{tab:related}
\centering
\begin{reporttable}
\begin{tabularx}{\textwidth}{|L{2.6cm}|Y|Y|}
\hline
\textbf{System} & \textbf{Who decides an update} & \textbf{Main limitation} \\
\hline
MemGPT & The model, by function calls & Edits depend on correct calls \\
\hline
Zep / Graphiti & Graph rule (validity intervals) & Single run; graph is the main store \\
\hline
Mem0 & The LLM, per item & One benchmark; code differs from paper \\
\hline
Supermemory & Not documented & Closed source, vendor report only \\
\hline
Hindsight & Retain step & Not peer reviewed \\
\hline
\end{tabularx}
\end{reporttable}
\end{table}

In each system the model decides updates or the rule is tied to a graph database
(Table~\ref{tab:related}). We kept extract-then-compare, validity intervals and fused recall.
Our design differs in four ways: the update rule is plain code, PostgreSQL is the source of
truth, every assertion has an exact quote, and a small local model does the extraction.
